\section*{Capítulo \ref{ch:b2:alpine-orogen} --- Anatomia de um cinturão de montanhas: os Alpes}

\begin{solution}{exo:b2:alpine-orogen:1}
(a) $e = 1 - 27.0/36.0 = 0.25$, isto é, \qty{25}{\percent}; $S = 36.0 - 27.0 =
\qty{9.0}{km}$. (b) A camada deve ter conservado seu comprimento (sem estiramento nem adelgaçamento ao longo dela),
a seção deve ser paralela à direção de transporte, e as linhas de fixação devem marcar lugares
onde as camadas não deslizaram umas sobre as outras.
\end{solution}

\begin{solution}{exo:b2:alpine-orogen:2}
(a) Uma \omterm{def:b2:alpine-orogen:klippe}{klippe}: o gnaisse é um resto isolado de uma \omterm{def:b2:alpine-orogen:nappes}{nappe} repousando sobre seu autóctone;
a erosão removeu o resto da lâmina em todo o seu redor. (b) Uma \omterm{def:b2:alpine-orogen:klippe}{janela tectônica}: a erosão
cortou a \omterm{def:b2:alpine-orogen:nappes}{nappe} de gnaisse e expôs as rochas mais jovens abaixo do seu cavalgamento basal, que
por isso mergulha para longe da janela em todos os lados.
\end{solution}

\begin{solution}{exo:b2:alpine-orogen:3}
(a) Complexo ultramáfico (o manto superior); complexo gabroico (a crosta oceânica inferior);
complexo de \omterm{def:b2:magma-emplacement:bodies}{diques} em lençol (a crosta média); complexo vulcânico, em almofadas (a crosta superior).
(b) O complexo de \omterm{def:b2:magma-emplacement:bodies}{diques} em lençol. (c) A maior parte da crosta dos oceanos largos é consumida pela subducção;
os \omterm{def:b2:alpine-orogen:ophiolite}{ofiolitos} que sobrevivem em terra se formaram sobretudo em pequenas bacias oceânicas ou entre arcos
insulares, e os tetianos onde um continente se afinou até que um pequeno oceano se abrisse.
\end{solution}

\begin{solution}{exo:b2:alpine-orogen:4}
(a) \omterm{def:b2:alpine-orogen:nappes}{Nappes} do Monte Rosa e do Grand St Bernard: o lado europeu (o Briançonnais e a
margem passiva europeia); Zermatt--Saas e zona do Combin: o oceano Piemonte--Ligúria; sistema da Dent
Blanche: a margem da Apúlia. (b) A Apúlia, que termina por cima. (c) A placa superior não é
subductada, e as lascas raspadas da placa inferior se empilham abaixo dela na ordem em que a
placa inferior chegou: o oceano distante primeiro, a margem próxima por último e mais embaixo. Um dobramento posterior da pilha,
ou cavalgamentos fora de sequência, podem perturbar a ordem.
\end{solution}

\begin{solution}{exo:b2:alpine-orogen:5}
(a) $2\pi G\,\Delta\rho = 2\pi\times6.674\times10^{-11}\times500 =
\qty{2.097e-7}{s^{-2}}$; $h = 1.70\times10^{-3}/2.097\times10^{-7} = \qty{8.1}{km}$.
(b) $2.097\times10^{-7}\times25\,000 = \qty{5.24e-3}{m.s^{-2}}$, cerca de
\qty{-524}{mGal}. (c) A raiz não é uma placa infinita: é um corpo de algumas dezenas de quilômetros de
profundidade e com mais ou menos a largura do cinturão, de modo que ela atrai menos um ponto da superfície do que o faria uma
placa da mesma espessura estendida ao infinito; corpos densos como a zona de Ivrea
acrescentam anomalias positivas, e o contraste de densidade em profundidade pode ser menor do que o suposto.
\end{solution}

\begin{solution}{exo:b2:alpine-orogen:6}
(a) $164 - 44 = \qty{120}{Myr}$. (b) $z = P/(\bar\rho g)$:
$2.8\times10^9/(3000\times9.8) = \qty{95.2}{km}$; $2.8\times10^9/(3300\times9.8) =
\qty{86.6}{km}$. (c) ``Mais de \qty{90}{km}'' fica dentro desse intervalo: a profundidade depende
tanto da densidade média escolhida para a sobrecarga quanto da pressão.
\end{solution}

\begin{solution}{exo:b2:alpine-orogen:7}
(a) Um flanco de comprimento $L$ mergulhando com o ângulo $\theta$ se estende por $L\cos\theta$ na horizontal; todos os flancos
têm o mesmo mergulho, logo $l = l_0\cos\theta$ e $e = 1 - l/l_0 = 1 - \cos\theta$.
(b) $\cos\theta = 0.70$: $\theta = \qty{45.6}{\degree}$. (c) $1 - \cos\qty{40}{\degree} =
\qty{23.4}{\percent}$, como na figura dos chevrons; $1 - \cos\qty{60}{\degree}
= \qty{50}{\percent}$.
\end{solution}

\begin{solution}{exo:b2:alpine-orogen:8}
(a) $l_0 = 118/2.0 = \qty{59.0}{km}$; $S = 59.0 - 42 = \qty{17.0}{km}$; $e = 17.0/59.0 =
\qty{28.8}{\percent}$. (b) $A_e = 118 - 42\times2.0 = \qty{34.0}{km^2}$; $d = 34.0/17.0 =
\qty{2.0}{km}$. (c) Sim: o \omterm{def:b2:alpine-orogen:decollement}{descolamento} fica na base da cobertura de \qty{2.0}{km}, abaixo
do seu topo; a cobertura está destacada sobre sua base.
\end{solution}

\begin{solution}{exo:b2:alpine-orogen:9}
(a) A raiz de Airy é $4.5h$: a \qty{46.5}{\degree} N, $4.5\times2.74 = \qty{12.3}{km}$, um
Moho a \qty{42.3}{km}; a \qty{45.55}{\degree} N, $4.5\times0.65 = \qty{2.9}{km}$, um Moho
a \qty{32.9}{km}. (b) Sob o terreno alto, o Moho está \qty{6.0}{km} mais raso do que Airy
prevê; sob a borda da planície do Pó, está \qty{28.3}{km} mais fundo. (c) A raiz não é
compensada localmente: ela está deslocada para o sul e é maior do que o relevo exige.
A placa se curva como um todo sob as cargas que pesam sobre ela (o suporte regional de um capítulo do terceiro
ano), e uma dessas cargas, o corpo de Ivrea, não é topografia de modo algum.
\end{solution}

\begin{solution}{exo:b2:alpine-orogen:10}
(a) $5.1/11.3 = 0.45$: menos da metade. (b) $5.1 + 6.2 = \qty{11.3}{km}$. O topo do
Jurássico fica acima do \omterm{def:b2:alpine-orogen:decollement}{descolamento} superior; os \qty{6.2}{km} transportados ao longo desse \omterm{def:b2:alpine-orogen:decollement}{descolamento}
passaram por baixo das dobras da seção e foram gastos em estruturas mais ao norte, de modo que a
camada superior nunca se encurtou desse valor dentro da seção. (c) Uma única camada-guia registra
apenas o encurtamento das camadas a que pertence; cada \omterm{def:b2:alpine-orogen:decollement}{descolamento} acima ou abaixo dela pode
levar deslocamento para dentro ou para fora da seção. Balanceiem-se várias camadas, e confira-se com as áreas.
\end{solution}

\begin{solution}{exo:b2:alpine-orogen:11}
(a) Para um rejeito de \qty{100}{km}, $\sqrt{200^2 + 100^2} = \qty{223.6}{km}$; para
\qty{50}{km}, $\sqrt{200^2 + 50^2} = \qty{206.2}{km}$; para \qty{150}{km},
$\sqrt{200^2 + 150^2} = \qty{250.0}{km}$. (b) $\arctan(100/200) = \qty{26.6}{\degree}$ em relação à
direção norte--sul, uma linha de convergência próxima da N\qty{150}{\degree}E citada por Rosenberg e
colaboradores. (c) O rejeito é incerto em $\pm\qty{50}{km}$, o que desloca a convergência
total em cerca de \qty{20}{km} para um lado ou para o outro; o terceiro algarismo do resultado está muito abaixo
dessa incerteza.
\end{solution}

\begin{solution}{exo:b2:alpine-orogen:12}
(a) $210/12.9 = \qty{16.3}{Myr}$ a $210/8.5 = \qty{24.7}{Myr}$ (\unit{km} dividido por
\unit{mm.yr^{-1}} dá milhões de anos). (b) Mais de $4\times60 = \qty{240}{km}$; pelo
menos $240 - 60 = \qty{180}{km}$ da bacia estão hoje sob os cavalgamentos alpinos ou foram
erodidos. (c) A bacia se forma na frente do prisma de cavalgamentos e migra com ele: sua parte
interna está sempre sendo cavalgada pelos cavalgamentos que a carregam, e seu preenchimento mais antigo fica preso
nos cavalgamentos, como a \omterm{def:b2:alpine-orogen:foreland}{Molassa} Subalpina.
\end{solution}

\begin{solution}{pb:b2:alpine-orogen:1}
\noindent\textbf{Parte I.}

\textbf{1.} $3.0/\cos\qty{35}{\degree} = \qty{3.66}{km}$; $2.5/\cos\qty{50}{\degree} =
\qty{3.89}{km}$; $4.0/\cos\qty{40}{\degree} = \qty{5.22}{km}$.

\textbf{2.} $8.00 + 2\times3.66 + 5.00 + 2\times3.89 + 6.0/\cos\qty{10}{\degree} +
2\times5.22 + 10.0/\cos\qty{5}{\degree} = 8.00 + 7.32 + 5.00 + 7.78 + 6.09 + 10.44 + 10.04 =
\qty{54.67}{km}$; com o trecho repetido, $l_0 = \qty{60.7}{km}$.

\textbf{3.} A camada sob a lâmina de cavalgamento é a mesma camada, levada para baixo da lâmina pelo
cavalgamento: antes do cavalgamento, ela estava alinhada com o resto. Deixá-la de fora restauraria a
camada curta demais.

\textbf{4.} $e = 1 - 48.0/60.7 = 0.209$, isto é, \qty{20.9}{\percent}; $S = 60.7 - 48.0 =
\qty{12.7}{km}$.

\textbf{5.} $60.7 - 60.2 = \qty{0.5}{km}$, $0.5/60.7 = \qty{0.8}{\percent}$: dentro da
precisão da medida; a seção é aceitável.

\textbf{6.} $60.7 - 54.0 = \qty{6.7}{km}$, \qty{11}{\percent}: as duas camadas de uma mesma lâmina
não podem se restaurar em comprimentos tão diferentes. A seção não seria admissível tal como traçada: uma
estrutura faltando ou mal desenhada, uma linha de fixação no lugar errado, movimento para fora do plano da
seção, ou um \omterm{def:b2:alpine-orogen:decollement}{descolamento} entre as duas camadas.

\noindent\textbf{Parte II.}

\textbf{7.} $146/2.4 = \qty{60.8}{km}$.

\textbf{8.} Ele difere de \qty{60.7}{km} em \qty{0.1}{km}: o comprimento e a área concordam.

\textbf{9.} $A_e = 146 - 48.0\times2.4 = 146 - 115.2 = \qty{30.8}{km^2}$.

\textbf{10.} $d = 30.8/12.7 = \qty{2.43}{km}$, igual aos \qty{2.4}{km} da cobertura dentro da
precisão dos dados: o \omterm{def:b2:alpine-orogen:decollement}{descolamento} é a base da cobertura.

\textbf{11.} $160/2.4 = \qty{66.7}{km}$, \qty{6.0}{km} a mais do que o comprimento da camada. O espessamento
dos evaporitos na base da cobertura, que fluem para os núcleos das anticlinais, acrescenta área
sem acrescentar comprimento; o balanço de área com espessura constante então falha.

\textbf{12.} A cobertura se moveu sobre um \omterm{def:b2:alpine-orogen:decollement}{descolamento} em sua base e o embasamento não foi
envolvido. Sob o cinturão, o topo do embasamento deve ser plano ou mergulhar suave e
regularmente, sem ser cortado pelos cavalgamentos acima dele.

\noindent\textbf{Parte III.}

\textbf{13.} $22 + 18 + 25 = \qty{65}{km}$; $S = 65 - 15 = \qty{50}{km}$.

\textbf{14.} $e = 50/65 = \qty{76.9}{\percent}$, contra \qty{20.9}{\percent} no
cinturão de antepaís: as \omterm{def:b2:alpine-orogen:nappes}{nappes} são muito mais encurtadas.

\textbf{15.} $30 - 22 = \qty{8}{km}$ ($e = \qty{26.7}{\percent}$). O próprio embasamento está
encurtado: \omterm{def:b2:alpine-orogen:skin}{tectônica de pele grossa}.

\textbf{16.} Cinturão de antepaís: $12.7/(11 - 4) = \qty{1.81}{km/Myr}$, isto é,
\qty{1.81}{mm} por ano; \omterm{def:b2:alpine-orogen:nappes}{nappes}: $50/(30 - 20) = \qty{5.0}{mm}$ por ano.

\textbf{17.} As \omterm{def:b2:alpine-orogen:nappes}{nappes} (30 a \qty{20}{Ma}) se deformaram antes do cinturão de antepaís (11 a
\qty{4}{Ma}): a deformação migrou para fora, em direção ao antepaís. Essa é a ordem dos
Alpes: o cavalgamento de Glarus no Oligoceno e no Mioceno, os maciços externos no Mioceno,
o Jura no Mioceno Superior e no Plioceno.

\textbf{18.} $12.7 + 50 + 8 = \qty{70.7}{km}$.

\noindent\textbf{Parte IV.}

\textbf{19.} $\qty{0.5}{cm.yr^{-1}}\times\qty{40}{Myr} = 0.005\times4\times10^7 =
\qty{2.0e5}{m} = \qty{200}{km}$.

\textbf{20.} $200 + 18 = \qty{218}{km}$.

\textbf{21.} Só se conta o que está preservado ou é inferido: partes do oceano podem ter
sido subductadas antes de \qty{85}{Ma} ou mais depressa do que a taxa média, lascas de \omterm{def:b2:alpine-orogen:ophiolite}{ofiolito} podem ter
sido erodidas ou levadas para fora da seção, e uma convergência oblíqua à seção traz
fundo oceânico que a seção não vê.

\textbf{22.} Zona interna: $60 - 20 = \qty{40}{km}$ ($e = \qty{66.7}{\percent}$); retrocunha:
$70 - 48 = \qty{22}{km}$ ($e = \qty{31.4}{\percent}$).

\textbf{23.} $12.7 + 50 + 8 + 40 + 22 = \qty{132.7}{km}$; dentro do intervalo das estimativas
para os Alpes Centrais, cerca de 125 a \qty{160}{km} de encurtamento colisional (antes de somar
o deslizamento direcional).

\textbf{24.} Largura restaurada: $153 + 132.7 = \qty{285.7}{km}$; crosta restaurada:
$285.7\times30 = \qty{8571}{km^2}$; crosta atual: $153\times45 = \qty{6885}{km^2}$. Falta:
$8571 - 6885 = \qty{1686}{km^2}$, \qty{19.7}{\percent} da crosta restaurada; ela foi
erodida (e despejada nas \omterm{def:b2:alpine-orogen:foreland}{bacias de antepaís}) ou subductada com a crosta inferior para dentro do
manto. Sem essas perdas, a crosta teria em média $8571/153 = \qty{56}{km}$ de espessura.
O encurtamento colisional total através da seção é de
\textbf{\qty[text-series-to-math=true]{133}{km}}.
\end{solution}
